\section{Rounds of Adaptivity}\label{app:roundshier}
\paragraph{Interviews with $r$ rounds.}
An $r$-round interview chooses coordinate batches
$B_1,\;B_2(x_{B_1}),\;\ldots,\;B_r(x_{B_1\cup\cdots\cup B_{r-1}})$.
Each batch is a function of the answers to earlier batches only, and
$\sum_s|B_s|\le b$ on every input. Let $\A_r(b)$ be the finite menu of
deterministic $r$-round interviews and $V_r(d,b)=V(d,\A_r(b))$; randomized
interviews are mixtures and are already represented by the design
probabilities in \eqref{eq:unbiasedprogram}. Fixed forms are $\A_1(b)$.
Sequentializing each batch shows that every $r$-round interview is a depth-$b$
tree, and a depth-$b$ tree is a $b$-round interview with one query per round.
Hence $\A_1(b)\subseteq\A_2(b)\subseteq\cdots\subseteq\A_b(b)=\A_{\rm branch}(b)$
and $V_1\ge V_2\ge\cdots\ge V_b=V_{\rm branch}$. All menus identify the same
contrasts, $\deg(d)\le b$, because the two extreme menus do
(Appendix~\ref{app:id}). Theorems~\ref{thm:finite} and~\ref{thm:asym} apply to
each finite menu $\A_r(b)$, so the lower bounds below hold against biased,
sequential and label-dependent procedures.

\begin{theorem}[Rounds of adaptivity, full statement of Theorem~\ref{thm:rounds}]\label{thm:roundsfull}
\textup{(a)} For $1\le r\le b\le m$ and every nonzero contrast of degree at most $b$ on $\{-1,1\}^m$,
\[
 V_{\rm branch}(d,b)\le V_r(d,b)\le
 \min\Big\{2^{b-r},\binom{m-r}{b-r}\Big\}\,V_{\rm branch}(d,b).
\]
\textup{(b)} For integers $r,\kappa\ge1$ put $D=2^\kappa$ and $b=r\kappa+1$.
The pointer-chasing contrast $N_{r,\kappa}$ defined below satisfies
\[
 V_{r+1}(N_{r,\kappa},b)=V_{\rm branch}(N_{r,\kappa},b)=1,\qquad
 \frac{D}{b}\le V_r(N_{r,\kappa},b)\le D.
\]
\end{theorem}

\begin{corollary}[Two rounds can be exponentially worse than branching]\label{cor:tworounds}
For every $b\ge3$,
\[
 \sup_{m,d}\frac{V_2(d,b)}{V_{\rm branch}(d,b)}\ \ge\ \frac{2^{\lfloor(b-1)/2\rfloor}}{b},
 \qquad
 \sup_{m,d}\frac{V_r(d,b)}{V_{r+1}(d,b)}\ \ge\ \frac{2^{\lfloor(b-1)/r\rfloor}}{b}
 \quad(1\le r<b).
\]
\end{corollary}
\begin{proof}
Take $\kappa=\lfloor(b-1)/r\rfloor\ge1$ and $N_{r,\kappa}$, which uses
$r\kappa+1\le b$ queries along every path. Branching and $(r+1)$-round
interviews still compute it exactly under the larger cap $b$, while the hit
bound in the proof of part (b) uses only the total number $b$ of queries.
For the first display use $r=2$ and $V_{\rm branch}\le V_3$.
\end{proof}

The constant separation $V_2\ge4/3$ of Corollary~\ref{cor:roundseparation}
is therefore not the end of the hierarchy. Part (a) and the corollary leave
the exponent for $2\le r<b$ between $\lfloor(b-1)/r\rfloor-\log_2b$ and $b-r$;
the exact frontier is open.

\paragraph{Proof of part (a).}
The left inequality is menu inclusion. For the right inequality pad a
depth-$b$ tree $\tau$ as in Appendix~\ref{app:universal}, so that every path
has $b$ distinct coordinates. In rounds $1,\ldots,r-1$ ask one question per
round and follow $\tau$ exactly. This reaches a node $v$ of depth $r-1$,
whose padded subtree has depth $b-r+1$. Its path sets all contain the subtree
root and are determined by the first $b-r$ answers below $v$; hence their
number $L_v$ is at most $2^{b-r}$, and also at most $\binom{m-r}{b-r}$
because the remaining $b-r$ coordinates are chosen from the $m-r$
coordinates not on the followed prefix or at $v$. In round $r$ choose one of
these $L_v$ sets uniformly, independently of the person, and attempt to
complete the path. On completion return $L_v$ times the tree score and
otherwise zero. Conditionally on $x$ the mean equals the tree score and the
second moment is $L_v$ times its square. The mixture and merging argument of
Appendix~\ref{app:universal} applies verbatim to randomized branching
designs, so every feasible branching design is simulated with second
moment inflated by at most $\max_vL_v$.

\paragraph{The pointer-chasing contrast.}
Levels $\ell=1,\ldots,r$ consist of registers. Level $\ell$ has $D^{\ell-1}$
registers $R_\ell(i_1,\ldots,i_{\ell-1})$, indexed by
$(i_1,\ldots,i_{\ell-1})\in[D]^{\ell-1}$, and each register stores $\kappa$
binary coordinates encoding an element of $[D]$. There are $D^r$ data
coordinates $Z_{i_1\cdots i_r}$. The pointer path is $j_1=R_1$ and
$j_\ell=R_\ell(j_1,\ldots,j_{\ell-1})$, and
\[
 N_{r,\kappa}(x)=Z_{j_1\cdots j_r}
 =\sum_{i\in[D]^r}Z_i\prod_{\ell=1}^r\mathbf 1\{R_\ell(i_{<\ell})=i_\ell\}.
\]
Each register indicator has degree $\kappa$, so $\deg N_{r,\kappa}=r\kappa+1=b$.
Reading the $\kappa$ coordinates of the level-$\ell$ register on the pointer
path in round $\ell$, and the selected data coordinate in round $r+1$,
computes the contrast with $b$ queries in $r+1$ rounds. With Jensen's
inequality this gives $V_{r+1}=V_{\rm branch}=1$. For the upper bound
$V_r\le D$, read the pointer registers of levels $1,\ldots,r-1$ in the first
$r-1$ rounds. In round $r$, read the level-$r$ register together with one
data coordinate $Z_{j_1\cdots j_{r-1}J}$, where $J$ is uniform on $[D]$ and
drawn in advance. The score $DZ_{j_1\cdots j_{r-1}J}\mathbf 1\{J=j_r\}$ is
unbiased with second moment $D$ at every profile.

\paragraph{The lower bound $V_r(N_{r,\kappa},b)\ge D/b$.}
Let $\U$ be the uniform law and write $c^\star=(j_1,\ldots,j_r)$ for the
selected data cell. Since $\E_\U[N^2]=1$, the dual formula~\eqref{eq:vdual}
with witness $h=N$ gives $V_r\ge1/\max_{\iota\in\A_r(b)}E_\iota$, where
$E_\iota=\E_\U[\E_\U[N\mid\ell_\iota(X)]^2]$. Fix a deterministic $r$-round
interview $\iota$.

\emph{Step 1: a generous transcript.}
Let $\iota^+$ ask the same questions as $\iota$, but at the end of every
round additionally reveal the full content of every register touched in that
round. The transcript of $\iota$ is a function of that of $\iota^+$, and
conditional expectation onto a larger $\sigma$-field cannot decrease the
second moment, so $E_\iota\le E_{\iota^+}$.

\emph{Step 2: only the selected data coordinate carries signal.}
Let ${\rm Obs}$ be the set of data coordinates queried by $\iota^+$. For a
data cell $c$, on the transcript event $\{c\notin{\rm Obs}\}$ the transcript
and the pointer path are functions of the coordinates other than $Z_c$,
because the interview never reads $Z_c$. Since $Z_c$ is an independent
uniform sign, $\E[Z_c\mathbf 1\{c^\star=c\}\mid\ell_{\iota^+}]$ vanishes there.
Hence
\[
 \E[N\mid\ell_{\iota^+}]=\sum_{c\in{\rm Obs}}Z_c\,\PP(c^\star=c\mid\ell_{\iota^+}),
\]
so $|\E[N\mid\ell_{\iota^+}]|\le\PP(c^\star\in{\rm Obs}\mid\ell_{\iota^+})\le1$ and
\[
 E_{\iota^+}\le\E\big[\PP(c^\star\in{\rm Obs}\mid\ell_{\iota^+})\big]
 =\PP(c^\star\in{\rm Obs}).
\]

\emph{Step 3: observing $c^\star$ in $r$ rounds requires a lucky guess.}
Let $\G_t$ be the $\sigma$-field of the $\iota^+$ transcript before round $t$,
and let $K_t$ be the largest $k$ such that the pointer registers
$R_1,R_2(j_1),\ldots,R_k(j_{<k})$ are all revealed before round $t$; thus
$K_1=0$ and $K_t$ is nondecreasing. A query of round $t$ is \emph{deep} if it
touches a register of level at least $K_t+2$ or a data coordinate while
$K_t\le r-1$; it is a \emph{hit} if, in addition, the touched register or
data coordinate lies on the pointer path. Every deep location is indexed by a
tuple whose $(K_t+1)$st entry is $\G_t$-measurable, and lies on the pointer
path only if this entry equals $j_{K_t+1}$. The latter is the content of the
register $R_{K_t+1}(j_{\le K_t})$, which is unrevealed before round $t$.
For an adaptive sequence of queries to independent uniform coordinates, the
unrevealed coordinates remain independent and uniform given the
transcript. Hence $j_{K_t+1}$ is uniform on $[D]$ given $\G_t$, and every deep
query is a hit with conditional probability at most $1/D$.

If no hit occurs in rounds $1,\ldots,t$, then $K_{t+1}\le K_t+1$. Indeed,
$K_{t+1}\ge K_t+2$ requires the pointer register of level $K_t+2$ to be
revealed by the end of round $t$, say in round $t'\le t$. Since
$K_{t'}\le K_t$, that query was deep in round $t'$ and lay on the pointer path,
so it was a hit. By induction, without hits $K_t\le t-1\le r-1$ for every
$t\le r$. Then every query of a data coordinate in rounds $1,\ldots,r$ is deep,
and $c^\star$ can be queried only by a hit. Since the total number of queries
is at most $b$,
\[
 \PP(c^\star\in{\rm Obs})\le\PP(\text{some hit})
 \le\E\Big[\sum_{t=1}^r\sum_{\text{deep queries of round }t}\PP(\text{hit}\mid\G_t)\Big]
 \le\frac bD .
\]
Combining the three steps gives $E_\iota\le b/D$ for every $\iota\in\A_r(b)$,
and therefore $V_r\ge D/b$. \hfill$\square$

\paragraph{A binary pointer tree: few rounds lose almost the whole exponent.}
The pointer-chasing family separates consecutive numbers of rounds, but its
lower bound $2^{\lfloor(b-1)/r\rfloor}/b$ leaves the exponent for a fixed
small number of rounds far below the upper bound $b-r$ of part~(a). The
following contrast closes most of this gap. Let $L\ge1$ and $b=L+1$. Arrange
$2^L-1$ \emph{pointer coordinates} as the internal nodes $u$ of a complete
binary tree with levels $0,\ldots,L-1$ and attach $2^L$ \emph{data
coordinates} $Z_v$ to its leaves $v$, so that $m=2^{b}-1$. Encode the pointer
coordinate of node $u$ as a bit $x_u\in\{0,1\}$ that selects one of the two
children of $u$. The \emph{pointer path} starts at the root and descends
through the selected children; let $c^\star(x)$ be the leaf it reaches, and
\[
 B_L(x)=Z_{c^\star(x)}=\sum_{v}Z_v\prod_{u\in{\rm path}(v)}\mathbf 1\{x_u\to v\},
\]
where ${\rm path}(v)$ is the set of the $L$ internal nodes above $v$ and
$\{x_u\to v\}$ is the event that the bit of $u$ selects the child on the way
to $v$. Each indicator has degree one, so $\deg B_L=L+1=b$, and following
the path with one question per round computes $B_L$ with $b$ queries. The
function is the tree evaluation problem of \citet{cook2012}; a deterministic
adversary argument on the same instance shows that exact evaluation with at
most $B$ queries per round needs $L/\lceil\log_2(B+1)\rceil$ rounds
\citep{li2026regex}. Here the total number of queries is capped at $b$, exact
evaluation in fewer than $b$ rounds is impossible for every strategy, and we
quantify the variance that the missing rounds cost.

\begin{theorem}[Binary pointer tree]\label{thm:binarytree}
Let $L\ge1$, $b=L+1$ and let $B_L$ be defined on $\{0,1\}^{2^b-1}$ as above.
\textup{(a)} $V_{\rm branch}(B_L,b)=V_b(B_L,b)=1$ and $V_r(B_L,b)\le2^{b-r}$
for $1\le r\le b$.
\textup{(b)} $V_1(B_L,b)\ge2^{b-1}/b$.
\textup{(c)} For $2\le r\le b$,
\[
 V_r(B_L,b)\ \ge\ \frac{2^{b-1}}{r\,b\,\big(1+\tfrac{2b}{r-1}\big)^{r-1}}
 \ \ge\ \frac{2^{b-1}}{(3b)^{r}} .
\]
\end{theorem}

\begin{corollary}[Few rounds lose almost the whole exponent]\label{cor:fewrounds}
For every $b\ge2$,
\[
 \frac{2^{b-2}}{b(2b+1)}\ \le\ \sup_{m,d}\frac{V_2(d,b)}{V_{\rm branch}(d,b)}\ \le\ 2^{b-2},
 \qquad
 \frac{2^{b-1}}{(3b)^{r}}\ \le\ \sup_{m,d}\frac{V_r(d,b)}{V_{\rm branch}(d,b)}\ \le\ 2^{b-r}
 \quad(2\le r\le b).
\]
Hence the exponent of the $r$-round gap lies between $b-1-r\log_2(3b)$ and
$b-r$: for every fixed $r$ it equals $b-O(r\log b)$, and an $r$-round protocol
that is within a factor $b^{c}$ of full branching on every contrast of degree
at most $b$ needs $r\ge(b-1-c\log_2b)/\log_2(3b)$ rounds.
\end{corollary}
\begin{proof}
The upper bounds are Theorem~\ref{thm:roundsfull}(a). The lower bounds are
Theorem~\ref{thm:binarytree}(c) with $V_{\rm branch}(B_L,b)=1$; for $r=2$
the first expression in~(c) equals $2^{b-1}/(2b(1+2b))$.
\end{proof}

The lower bounds are informative only once $2^{b-2}>b(2b+1)$, that is for
$b\ge10$; for smaller $b$ the constant separation $V_2\ge\tfrac43V_{\rm
branch}$ of Corollary~\ref{cor:roundseparation} at $b=3$ remains the best
statement we have. Like the address contrast, $B_L$ uses a number of
coordinates exponential in $b$. Whether the polynomial factors in~(c) can be
removed, so that the $r$-round gap is $2^{b-r}$ up to constants, is open.

\paragraph{Proof of Theorem~\ref{thm:binarytree}.}
(a) Following the pointer path returns $B_L$ itself, so $V_b\le1$; every
unbiased design has $\sum_\iota g_\iota(x)=B_L(x)=\pm1$, and Cauchy--Schwarz
gives $\sum_\iota g_\iota(x)^2/q_\iota\ge1$ at every $x$, so $V_{\rm branch}\ge1$.
The remaining inequality is Theorem~\ref{thm:roundsfull}(a).

(b) and (c). Let $\U$ be the uniform law on all coordinates, under which
$\E_\U[B_L^2]=1$ and the pointer path is a uniformly distributed root-to-leaf
path. By the dual formula~\eqref{eq:vdual} with witness $h=B_L$,
$V_r(B_L,b)\ge1/\max_{\iota\in\A_r(b)}E_\iota$ with
$E_\iota=\E_\U[\E_\U[B_L\mid\ell_\iota(X)]^2]$. Fix a deterministic $r$-round
interview $\iota$. Let $S_t$ be the set of coordinates it reads in round $t$;
$S_t$ is a function of the transcript $T_{t-1}$ of rounds $1,\ldots,t-1$, and
we may assume that no coordinate is read twice. Write $a_t=|S_t\cap U|$ and
$\lambda_t=|S_t\cap V|$ for the numbers of pointer and data coordinates read
in round $t$, so that $\sum_t(a_t+\lambda_t)\le b$, and let $I_t$ be the set of
pointer coordinates read in rounds $1,\ldots,t$ and ${\rm Obs}$ the set of
leaves whose data coordinate is read at all.

\emph{Step 1: only the selected leaf carries signal.}
On the event $\{c^\star\notin{\rm Obs}\}$, the transcript and the pointer
path are functions of coordinates other than $Z_{c^\star}$, so
$\E_\U[Z_v\mathbf 1\{c^\star=v\}\mid\ell_\iota]=0$ for every unread leaf $v$ and
\[
 \E_\U[B_L\mid\ell_\iota]=\sum_{v\in{\rm Obs}}Z_v\,\PP(c^\star=v\mid\ell_\iota),
 \qquad
 E_\iota\le\E\big[\PP(c^\star\in{\rm Obs}\mid\ell_\iota)^2\big]\le\PP(c^\star\in{\rm Obs}).
\]

\emph{Step 2: the conditional law of the selected leaf.}
For a leaf $v$ let $R_t(v)=|{\rm path}(v)\cap I_t|$ be the number of pointer
coordinates on its path read by the end of round $t$, and let $C_t(v)$ be the
event that every one of them selects the child towards $v$. The event
$\{T_t=\tau\}$ fixes the values of the coordinates read under $\tau$ and
nothing else, so conditionally on $T_t$ the unread coordinates remain
independent and uniform. Since $c^\star=v$ exactly when all $L$ pointer
coordinates on ${\rm path}(v)$ select the child towards $v$,
\begin{equation}
 \PP(c^\star=v\mid T_t)=\mathbf 1\{C_t(v)\}\,2^{R_t(v)-L}.\label{eq:leaflaw}
\end{equation}
For $r=1$ this gives $\PP(c^\star\in{\rm Obs})\le\lambda_1 2^{-L}\le b2^{-L}$,
which with Step~1 proves~(b).

\emph{Step 3: a potential for the most likely leaf.}
Put $M_t=\max_v\mathbf 1\{C_t(v)\}2^{R_t(v)}$, so that $M_0=1$ and, by
\eqref{eq:leaflaw}, $\max_v\PP(c^\star=v\mid T_t)=M_t2^{-L}$. We claim
\begin{equation}
 \E[M_t\mid T_{t-1}]\le(1+2a_t)\,M_{t-1}.\label{eq:potential}
\end{equation}
Condition on $T_{t-1}$, which fixes $S_t$. Let $v^\star$ attain $M_t$. If no
coordinate of $S_t$ lies on ${\rm path}(v^\star)$, then
$R_t(v^\star)=R_{t-1}(v^\star)$ and $C_t(v^\star)=C_{t-1}(v^\star)$, so
$M_t\le M_{t-1}$. Otherwise let $u$ be the deepest node of
${\rm path}(v^\star)\cap S_t$, write ${\rm path}(u)$ for the nodes above $u$
together with $u$ itself, and define
\[
 Q(u)=\mathbf 1\{\text{every node of }({\rm path}(u)\setminus\{u\})\cap I_t\text{ selects towards }u\}\,2^{|{\rm path}(u)\cap I_t|},
\]
\[
 B(u)=\max_{w\text{ below }u}\mathbf 1\{\text{every node of }({\rm path}(w)\setminus{\rm path}(u))\cap I_{t-1}\text{ selects towards }w\}\,2^{|({\rm path}(w)\setminus{\rm path}(u))\cap I_{t-1}|}.
\]
No node of ${\rm path}(v^\star)$ strictly below $u$ is read in round $t$, so
every read node there belongs to $I_{t-1}$ and
$\mathbf 1\{C_t(v^\star)\}2^{R_t(v^\star)}\le Q(u)B(u)$: the
factor $Q(u)$ accounts for the nodes of ${\rm path}(u)$ other than $u$ and for
the count including $u$, the bound $\mathbf 1\{x_u\to v^\star\}\le1$ drops
$u$'s own selection, and $B(u)$ dominates the contribution of the nodes below
$u$ because $w=v^\star$ is admissible in its maximum. Consequently
$M_t\le M_{t-1}+\sum_{u\in S_t\cap U}Q(u)B(u)$. The quantity $B(u)$ is
$T_{t-1}$-measurable. Given $T_{t-1}$, the $|{\rm path}(u)\cap S_t|-1$ newly
read nodes strictly above $u$ select towards $u$ independently with
probability $1/2$ each, so
\[
 \E[Q(u)\mid T_{t-1}]
 =\mathbf 1\{\text{every node of }{\rm path}(u)\cap I_{t-1}\text{ selects towards }u\}\,
   2^{|{\rm path}(u)\cap I_{t-1}|}\cdot2 .
\]
Multiplying by $B(u)$ reassembles, for the leaf $w$ below $u$ attaining
$B(u)$, exactly $2\cdot\mathbf 1\{C_{t-1}(w)\}2^{R_{t-1}(w)}\le2M_{t-1}$, because
$u\notin I_{t-1}$. Summing over the $a_t$ nodes of $S_t\cap U$
proves~\eqref{eq:potential}.

\emph{Step 4: reading the selected leaf.}
Since $S_t$ is $T_{t-1}$-measurable, \eqref{eq:leaflaw} gives
$\PP(c^\star\in S_t\cap V\mid T_{t-1})\le\lambda_t\,M_{t-1}2^{-L}$.

\emph{Step 5: combination.}
Let $\Psi_t=M_t\prod_{s\le t}(1+2a_s)^{-1}$. As $a_t$ is $T_{t-1}$-measurable,
\eqref{eq:potential} says $\E[\Psi_t\mid T_{t-1}]\le\Psi_{t-1}$, hence
$\E[\Psi_t]\le\Psi_0=1$. For $2\le t\le r$ the pathwise constraint
$\sum_{s<t}a_s\le b$ and the inequality between arithmetic and geometric means
give $\lambda_t\prod_{s<t}(1+2a_s)\le b\,(1+2b/(t-1))^{t-1}\le b\,(1+2b/(r-1))^{r-1}$,
because $k\mapsto(1+2b/k)^k$ is increasing; the empty product at $t=1$ is
$1$. Therefore
\[
 \PP(c^\star\in{\rm Obs})=\sum_{t=1}^r\E\big[\PP(c^\star\in S_t\cap V\mid T_{t-1})\big]
 \le2^{-L}\sum_{t=1}^r\E\Big[\lambda_t\prod_{s<t}(1+2a_s)\,\Psi_{t-1}\Big]
 \le2^{-L}\,r\,b\,\Big(1+\frac{2b}{r-1}\Big)^{r-1}.
\]
With Step~1 and the dual formula this proves the first inequality in~(c). For
the second, $1+2b/(r-1)\le3b/(r-1)$ because $r-1\le b$, so the denominator is
at most $rb(3b)^{r-1}/(r-1)^{r-1}\le(3b)^r$, since $r\le3(r-1)^{r-1}$ for
$r\ge2$. \hfill$\square$

The argument is a potential-function version of the lucky-hit bound: a
round can raise the conditional probability of the most likely leaf only by
a factor linear in the number of pointer coordinates it reads, whereas
following the path raises it by a factor two per question. As a numerical check, the released code simulates the recursion: on trees of
depth $3$ and $4$ with random previously read sets and random new batches, the
largest observed ratio of $\E[M_t\mid T_{t-1}]$ to $(1+2a_t)M_{t-1}$ was
$0.73$; for greedy two- and three-round strategies with total budget $b$ on
trees of depth $6$, $8$ and $10$, the sample mean of $\Psi_r$ was between
$0.32$ and $0.57$, and the hit probabilities ($0.25$, $0.09$ and $0.04$ for two
rounds) lay far below the bound of Step~5, as they must.


\paragraph{Interpretation.}
Each round can advance the known part of the pointer path by at most one
level unless a question lands on an unknown pointer target, which happens
with probability at most $1/D$ per question. Additional questions within a
round therefore do not substitute for a round. Practical counterparts of
rounds are questionnaire pages, clinic visits and batches of laboratory
orders. The address contrast of Theorem~\ref{thm:universal} is the case
$r=1$, where two rounds already attain the branching coefficient. The
pointer-chasing family shows that this is not a general property of
two-stage designs.

\section{The Value of Adaptivity at a Given Population}\label{app:population}
For a menu $\A$ and a law $P$ on $\X\times\{-1,1\}$ write
$M_\A(P)=W_\A(P)+\Delta(P)^2$, the smallest second moment of a score that is
unbiased for $t(x,y)=yd(x)$ at every profile and label. At a full-support
law, $W_\A(P)$ is also the local asymptotic minimax risk over arbitrary
estimators (Proposition~\ref{prop:local}). For a deterministic depth-$b$ tree
$\tau$ that computes $d$, let $S_\tau(x)$ be its padded path set and
$\mathcal S_\tau$ the collection of these sets
(Appendix~\ref{app:universal}).

\begin{theorem}[Population-specific bracket, full statement of Theorem~\ref{thm:bracket}]\label{thm:bracketfull}
\textup{(a)} For every menu $\A$, law $P$ and function $h$ on
$\X\times\{-1,1\}$ with $\E_P[th]\ne0$,
\[
 M_\A(P)\ \ge\ \frac{\E_P[th]^2}{\max_{\iota\in\A}\E_P\big[\E_P[h\mid\ell_\iota(X),Y]^2\big]}.
\]
In particular $M_\A(P)\ge\E_P[t^2]/\Gamma_\A(P)$, where
$\Gamma_\A(P)=\max_\iota\E_P[\E_P[t\mid\ell_\iota(X),Y]^2]/\E_P[t^2]$ is the
largest fraction of the target's energy explained by a single interview.

\textup{(b)} If a depth-$b$ tree $\tau$ computes $d$, then
$M_{\rm branch}(P)=\E_P[t^2]$ and
\[
 \frac1{\Gamma_{\rm fixed}(P)}
 \ \le\ \frac{M_{\rm fixed}(P)}{M_{\rm branch}(P)}
 \ \le\ N_P(\tau):=\frac{\Big(\sum_{S\in\mathcal S_\tau}\E_P\big[t^2\mathbf 1\{S_\tau(X)=S\}\big]^{1/2}\Big)^2}{\E_P[t^2]}
 \ \le\ |\mathcal S_\tau|\le K_{m,b}.
\]

\textup{(c)} Both inequalities in (b) are equalities, with value $D$, for the
address contrast with $\kappa=b-1$ under the uniform law with independent
uniform labels.
\end{theorem}

The quantity $N_P(\tau)=\exp H_{1/2}(w)$ is the exponential of the R\'enyi
entropy of order $1/2$ of the weights
$w_S=\E_P[t^21\{S_\tau=S\}]/\E_P[t^2]$: it is an effective number of
question sets that the tree actually uses at $P$. Adaptivity can be valuable
only when the tree's questions diverge across many well-populated paths. It
is valuable when, in addition, no fixed set of $b$ questions explains most of
the target, so that $\Gamma_{\rm fixed}(P)$ is small; correlated features and
shared coordinates raise $\Gamma_{\rm fixed}$. The variance ratio is
recovered as $\rho_P=(M_{\rm fixed}-\Delta^2)/(\E_Pt^2-\Delta^2)$.

\begin{proof}
(a) Let $H$ be any feasible score with interview probabilities $q_\iota$ and
transcript tables $g_\iota$, so $\sum_\iota g_\iota(\ell_\iota(x),y)=t(x,y)$
everywhere and $\E[H^2]=\sum_\iota\|g_\iota\|_P^2/q_\iota\ge(\sum_\iota\|g_\iota\|_P)^2$
by Cauchy--Schwarz in $q$. For every $h$,
\[
 \E_P[th]=\sum_\iota\E_P\big[g_\iota\,\E_P[h\mid\ell_\iota(X),Y]\big]
 \le\Big(\sum_\iota\|g_\iota\|_P\Big)\max_\iota\big\|\E_P[h\mid\ell_\iota(X),Y]\big\|_P,
\]
because $g_\iota$ is a function of $(\ell_\iota(X),Y)$. Squaring and taking
the infimum over feasible scores proves the display; $h=t$ gives the second
statement.

(b) Following $\tau$ returns $t$ itself, and Jensen's inequality gives
$\E[H^2]\ge\E[\E[H\mid X,Y]^2]=\E_P[t^2]$ for every unbiased score. The lower
bound is (a) with $\A$ the fixed forms. For the upper bound choose
$q_S>0$ on $\mathcal S_\tau$, draw $S$ before the person answers, and return
$t(X,Y)/q_S$ if the tree's path completes inside $S$ and zero otherwise.
Completion occurs exactly when $S=S_\tau(X)$
(Appendix~\ref{app:universal}), so the score is unbiased at every profile and
label, and its second moment is
$\sum_Sq_S^{-1}\E_P[t^2\mathbf 1\{S_\tau(X)=S\}]$. Minimizing over the simplex
gives $q_S\propto\E_P[t^21\{S_\tau=S\}]^{1/2}$ and the value
$N_P(\tau)\E_P[t^2]$. Sets of zero weight are handled by letting their
probability tend to zero. The last inequalities are the power-mean
inequality and $|\mathcal S_\tau|\le K_{m,b}$.

(c) Under the uniform law the weights are equal over the $D$ path sets, so
$N_P=D$. With independent labels,
$\E[t\mid X_T,Y]=Y\E[d\mid X_T]$, whose energy is at most $2^{-\kappa}=1/D$
by~\eqref{eq:addressinformation}. Hence $1/\Gamma_{\rm fixed}\ge D$, and the
universal bound $M_{\rm fixed}\le D$ forces equality throughout.
\end{proof}

\paragraph{The bracket on the variance scale.}
Section~\ref{sec:experiments} reports the number of records a design needs
relative to branching, that is the ratio of variances
$W_\A(P)=M_\A(P)-\Delta(P)^2$ of pointwise-unbiased designs rather than the
ratio of second moments. The two scales differ when $\Delta(P)\neq0$, and the
bracket transfers as follows.

\begin{corollary}[Bracket on the variance scale]\label{cor:bracketvar}
Under the assumptions of Theorem~\ref{thm:bracketfull}(b), with
$\sigma_P^2=\E_P[t^2]-\Delta(P)^2=W_{\rm branch}(P)>0$,
\[
 \frac{\E_P[t^2]/\Gamma_{\rm fixed}(P)-\Delta(P)^2}{\sigma_P^2}
 \ \le\ \frac{W_{\rm fixed}(P)}{W_{\rm branch}(P)}
 \ \le\ \frac{N_P(\tau)\,\E_P[t^2]-\Delta(P)^2}{\sigma_P^2}.
\]
When the second-moment endpoints exceed one, the variance-scale endpoints are
larger than them, so the variance ratio is bracketed by a wider interval
than the second-moment ratio.
\end{corollary}
\begin{proof}
Every design in \eqref{eq:mainW} is unbiased for $\Delta(P)$, so its variance
is its second moment minus $\Delta(P)^2$; in particular
$W_{\rm fixed}=M_{\rm fixed}-\Delta^2$ and $W_{\rm branch}=\E_P[t^2]-\Delta^2$.
Subtract $\Delta^2$ from the three terms of Theorem~\ref{thm:bracketfull}(b)
multiplied by $\E_P[t^2]$ and divide by $\sigma_P^2$. The last statement is
the inequality $(a-c)/(b-c)\ge a/b$ for $a\ge b>c\ge0$.
\end{proof}
The tables of Section~\ref{sec:experiments} and Appendix~\ref{app:screening}
report this variance-scale bracket, and "inside the bracket" refers to it.


\paragraph{Contrasts that branching cannot compute.}
For an ensemble whose contrast is not computed by one depth-$b$ tree, part (a)
applies to both menus. Every depth-$b$ branching interview is a coarsening of
a complete depth-$b$ tree, and the energy of a tree is a sum over its leaves.
Hence $\max_\iota\E_P[\E_P[h\mid\ell_\iota,Y]^2]$ over all depth-$b$ trees is
computed exactly by dynamic programming over partial assignments, in time
$O(m^b2^b)$. This yields certified lower bounds on $M_{\rm branch}(P)$ for every
dual witness $h$. Appendix~\ref{app:forestexact} uses these bounds inside
column generation to compute the branching coefficient over all depth-three
trees.
